\DocumentMetadata{
  pdfversion=2.0,
  pdfstandard=ua-2,
  lang=en-US,
  tagging=on,
  tagging-setup={math/setup=mathml-SE,tabsorder=structure}
}
\documentclass[11pt,letterpaper]{article}
\usepackage[margin=0.68in,includefoot]{geometry}
\usepackage{newcomputermodern}
\usepackage{amsmath,amscd,mathtools}
\usepackage{tikz,adjustbox}
\providecommand{\square}{\mdlgwhtsquare}
\usepackage[hidelinks]{hyperref}
\hypersetup{
  pdftitle={Numerical Analysis First-Year Exam, May 2018},
  pdfauthor={University of Florida Department of Mathematics},
  pdfsubject={Graduate mathematics examination},
  pdfdisplaydoctitle=true,
  bookmarksopen=true
}
\setcounter{secnumdepth}{0}
\setlength{\parindent}{0pt}
\setlength{\parskip}{0.28em}
\setlength{\emergencystretch}{3em}
\setlength{\leftmargini}{2.15em}
\setlength{\leftmarginii}{2.65em}
\setlength{\leftmarginiii}{3em}
\renewcommand{\labelenumi}{\textbf{\theenumi.}}
\pagestyle{empty}
\raggedbottom
\allowdisplaybreaks
\newenvironment{examproblems}[1]{%
  \begin{enumerate}%
  \setcounter{enumi}{#1}%
  \setlength{\topsep}{0.35em}%
  \setlength{\partopsep}{0pt}%
  \setlength{\itemsep}{0.48em}%
  \setlength{\parsep}{0.18em}%
}{\end{enumerate}}
\newenvironment{examsubparts}{%
  \begin{enumerate}%
  \setlength{\topsep}{0.18em}%
  \setlength{\partopsep}{0pt}%
  \setlength{\itemsep}{0.16em}%
  \setlength{\parsep}{0.1em}%
}{\end{enumerate}}
\newenvironment{examinstructions}{%
  \par\begingroup\itshape
}{%
  \par\endgroup\vspace{0.18em}
}
\makeatletter
\renewcommand\section{\@startsection{section}{1}{0pt}%
  {-1.25ex plus -.25ex minus -.1ex}%
  {0.55ex plus .1ex}%
  {\normalfont\large\bfseries}}
\renewcommand{\@maketitle}{%
  \newpage\null\vskip -1em%
  {\footnotesize\bfseries
    \href{https://www.ufl.edu/}{University of Florida}, \href{https://math.ufl.edu/}{Department of Mathematics}\par}%
  \vskip -0.5em%
  \tagstructbegin{tag=Title}%
    {\LARGE\bfseries\href{https://gma.math.ufl.edu/exams/numerical-analysis-first-year/}{Numerical Analysis First-Year Exam}, May 2018\par}%
  \tagstructend%
  \vskip 0.25em%
  \tagmcbegin{artifact}%
    \hrule height 1pt%
  \tagmcend%
  \vskip 1em%
}
\makeatother
\title{Numerical Analysis First-Year Exam, May 2018}
\author{}
\date{}
\begin{document}
\maketitle
\thispagestyle{empty}
\begin{examinstructions}
Do all five (5) problems.
\end{examinstructions}
\begin{examproblems}{0}
\item
Consider 
\[
y' = 4ty;\quad t \in [0,1];\quad y(0)=2, \tag{1}
\]
 which has solution \(Y(t)=2e^{2t^2}\).
\begin{examsubparts}
\item[\textnormal{(a)}]
Using the Euler method error estimate find the value of~\(h\) required to ensure that the Euler method solution \(w_i\) of (1) satisfies \(|Y(t_i)-w_i|\leq 0.1\) for all~\(i\).
\item[\textnormal{(b)}]
Derive the Taylor method of order 2 for (1).
\end{examsubparts}
\item
Let \(g\in C^2([a,b])\) and \(p\in(a,b)\) with \(g(p)=p\), \(g'(p)=0\), \(g''(p)\neq 0\).
\begin{examsubparts}
\item[\textnormal{(a)}]
Show there is an \(\epsilon>0\) so that for all \(x\in[p-\epsilon,p+\epsilon]\), we have \(g^n(x)\to p\) as \(n\to\infty\).
\item[\textnormal{(b)}]
With~\(\epsilon\) as in part (a), show that for all \(x\in[p-\epsilon,p+\epsilon]\), we have \(|g(x)-p|\leq M|x-p|^2\), where \(M=\max\{|g''(x)|:|x-p|\leq\epsilon\}/2\).
\end{examsubparts}
\item
Find the second order (i.e. quadratic) least squares approximation to the function \(f(x)=x^4\) on the interval \([-1,1]\) with respect to the weight \(w(x)\equiv 1\) using the fact that the first three orthonormal polynomials on \([-1,1]\) with respect to this weight are 
\[
\varphi_0(x)=\sqrt{1/2},\quad \varphi_1(x)=\sqrt{3/2}\,x,\quad\text{and}\quad \varphi_2(x)=\sqrt{5/8}(3x^2-1).
\]
\item
This problem has two parts.
\begin{examsubparts}
\item[\textnormal{(a)}]
Assume \(N(h)\) is the computed approximation for~\(M\) for each \(h>0\) and 
\[
M=N(h)+C_1h+C_2h^2+C_3h^3+\cdots.
\]
 Use the values \(N(h)\), \(N(h/3)\), and \(N(h/9)\) to produce a \(O(h^3)\) approximation to~\(M\).
\item[\textnormal{(b)}]
Taylor’s formula yields the following. 
\[
f'(x_0)=\frac{1}{h}\bigl(f(x_0+h)-f(x_0)\bigr)-\frac{h}{2}f''(x_0)-\frac{h^2}{6}f'''(x_0)+O(h^3).
\]
 Use this with the extrapolation of part (a) to derive an \(O(h^3)\) formula for \(f'(x_0)\).
\end{examsubparts}
\item
Find \(a,b,c\) so that the quadrature formula 
\[
\int_0^2 f(x)\,dx=af(0)+bf(1)+cf(2)
\]
 has degree of precision as large as possible, and show it is no larger than your answer.
\end{examproblems}
\end{document}
